
Von Altermagneten spricht man, wenn ein Material trotz vollständig kompensierter magnetischer Ordnung impulsabhängige spinaufgespaltene Bänder aufweist.
In dieser Arbeit untersuchen wir Altermagneten unter dem Einfluss verschiedener magnetischer Texturen, von breathing-Verzerrungen des Gitters und von Rashba-Spin-Bahn-Kopplung (Rashba-SOC).\\

Dazu wird ein zweidimensionales Tight-Binding-Modell für die magnetische Ordnungen der T1- und T2-Phasen von $\ce{Mn3Pt}$, sowie für die von \ce{Mn3Ge} und \ce{Mn3Sn} eingeführt. Mithilfe einer unitären Transformation des Tight-Binding-Hamiltonians werden die Positionen innerhalb der Einheitszelle berücksichtigt, um zell-periodische Zustände zu erhalten. Zur Berechnung der Berry curvature werden die Fukui-Hatsugai-Suzuki-Methode sowie zusätzlich eine zusammengesetzte Darstellung für isolierte Bandunterräume, die Entartungen erhalten, verwendet.\\
Mithilfe des Hamiltonians werden anschließend die Bandstruktur, Fermi-Oberflächen-Spintextur sowie die Berry curvature und Chern-Zahl numerisch für die unterschiedlichen magnetischen Texturen berechnet.\\
Dabei werden Symmetriebedingungen für die magnetischen Texturen hergeleitet, welche schließlich mit den numerischen Ergebnissen verglichen werden.\\

Die T1-und T2-Phasen von \ce{Mn3Pt} zeigen äquivalente Ergebnisse in Abwesenheit von SOC in Berücksichtigung einer globalen Spinrotation. Analog dazu liefern \ce{Mn3Ge} und \ce{Mn3Sn} dieselben Ergebnisse, da diese ebenfalls mit einer globalen Spinrotation zueinander verwandt sind.\\
Zudem wird in \ce{Mn3Ge} und \ce{Mn3Sn} eine Symmetrie gefunden, welche ein $C_{3z}$-Verhalten aufweist in Abwesenheit von SOC, wobei der Spinraum invers gedreht wird. Dagegen zeigt \ce{Mn3Pt} ebenfalls ein $C_{3z}$-Verhalten auf, jedoch nun in sowohl Spin- als auch Ortsraum äquivalent, was durch eine Symmetrie auferlegt wird, welche auch in Anwesenheit von SOC hält.\\
Die numerischen Ergebnisse zeigen ein Verhalten, das für die T1-Phase von \ce{Mn3Pt} und für \ce{Mn3Ge} mit den gefundenen $\mathcal{TM}$-Ebenen, sowie für die T2-Phase von \ce{Mn3Pt} und für \ce{Mn3Sn} mit den gefundenen $\mathcal{M}$-Ebenen übereinstimmt.\\
Im Grenzfall keiner Breathing-Verzerrungen und ohne SOC, erzwingen Symmetrien eine verschwindende Berry curvature.
Es wird gezeigt, dass Breathing-Verzerrungen Symmetrien brechen, welche gerades Verhalten erzwingen, während die in Abwesenheit von SOC erhaltene Symmetrie eine im Impulsraum ungerade Berry curvature erlaubt und damit eine verschwindende Chern-Zahl erzwingt.\\

Mithilfe der hergeleiteten Symmetriebedingungen wird gezeigt, dass \ce{Mn3Pt} keine Nettomagnetisierung in der Ebene aufweisen kann, während eine solche von \ce{Mn3Ge} und \ce{Mn3Sn} in Anwesenheit von SOC entlang verschiedener von Symmetrien ausgewiesenen Richtungen möglich ist. 
Die untersuchten Materialien bleiben auch in Abwesenheit von SOC spinaufgespalten, wodurch sie als starke Altermagneten identifiziert werden. Aufgrund der erlaubten Nettomagnetisierung in der Ebene werden \ce{Mn3Ge} und \ce{Mn3Sn} als Altermagneten vom M-Typ eigeordnet, während die T1- und T2-Phasen von \ce{Mn3Pt} als Altermagnete vom S-Typ eingeordnet werden.